\section{Finite-bit construction and a uniform enclosure certificate}
\label{sec:finite-bit}

The nearly linear bound in Theorem~\ref{thm:main} concerns exact real
arithmetic. We now give a separate construction in a finite-information
model, with an explicit word-size bound and nearly quadratic cost.
Its geometric input is the small-gap implication already proved in
Lemma~\ref{lem:defect}; its value oracle uses safeguarded integer
inertia counts rather than Graeffe transformations.

\subsection{Input model and theorem}
\label{subsec:finite-model}

Fix integers $n\geq2$ and $d\geq2$, and put $\delta=2^{-d}$.
All minima in the following definitions are over nonnegative integers:
\begin{align}
 a&=\min\{k:2^k\geq4n/\delta\},
 &\beta&=2^{-a},\label{eq:bit-a}\\
 K&=\min\{k:2^k\geq16n^3\},
 &\ell&=\beta4^{-K},\label{eq:bit-K}\\
 m&=\min\{k:2^k\geq n^2\},
 &p&=a+2K+20+m,\qquad Q=2^p,\qquad\eta=Q^{-1}.
 \label{eq:bit-precision}
\end{align}
Here $Q$ denotes the grid denominator only in this section.
Let $T$ be a complex Hermitian tridiagonal matrix with $\norm{T}\leq1$.
The finite input is $n,d$ and a Hermitian tridiagonal proxy $T_0$:
each real diagonal entry and each real and imaginary off-diagonal
component of $T_0$ is an integer multiple of $\eta$, and differs from
the corresponding component of $T$ by at most $\eta$.
Only these finite data are given to the algorithm. The entries of
$T$ need not be rational or separately available to it.

This is an input-enclosure model, not an encoding of arbitrary exact
real numbers as finite strings. Under the norm promise the input has
$O(np)$ bits. For rational data, reading and rounding the original
representation has additional cost, separate from the construction
bound below.

\begin{theorem}[Finite-bit diagonal regularization]
\label{thm:finite-bit}
For every $n,d,T,T_0$ satisfying the preceding hypotheses, a
deterministic integer algorithm, given $n,d,T_0$, returns a dyadic
number $c$ and the real diagonal matrix $D=cX$, where
$X=\diag(1,\ldots,n)$, such that
\begin{equation}\label{eq:bit-guarantee}
 0\leq nc\leq\delta/4,\qquad
 \gap(T+D)>2^{-100}\frac{\delta^4}{n^{42}}.
\end{equation}
It uses
\begin{equation}\label{eq:bit-cost}
 O\bigl(n^2p(1+\log n)\bigr)
 \quad\text{integer operations on $O(p)$-bit integers},
\end{equation}
and hence $O(n^2p^3(1+\log n))$ elementary bit operations with
schoolbook arithmetic, with $O(np)$ bits of workspace. Here
$p=O(\log(n/\delta))$, and all implicit constants are universal.
There are no square-root, eigenvalue, unbounded-precision-real or
exact-real-sign primitives.

The algorithm also returns $O(np)$ bits of certificate data:
three sets of rational ordered-eigenvalue enclosures for
$T_0+(c-r)X$, $T_0+cX$ and $T_0+(c+r)X$, where $r=\ell/4$,
and an explicit rational convex-defect bound. An independent exact
rational verifier checks these enclosures and the defect inequality.
The same output and certificate prove~\eqref{eq:bit-guarantee}
simultaneously for every normalized Hermitian tridiagonal matrix in
the componentwise $\eta$-enclosure centered at $T_0$.
The verifier may use longer exact integers; the bounded-word
construction cost does not assert a bounded-word verification cost.
\end{theorem}

We prove the theorem in the following subsections. No novelty is
claimed for safeguarded Sturm counting or bisection. Their role here
is to give an explicit finite-bit implementation of convex-defect
selection together with a certificate for an entire input enclosure.

\subsection{Safeguarded integer inertia counts}
\label{subsec:bit-inertia}

Let $A$ be a Hermitian tridiagonal matrix whose real diagonal and
real and imaginary edge components lie on the grid $Q^{-1}\mathbb Z$.
Write its diagonal entries as $A_i/Q$ and its complex edge entries
as $(B_i+\mathrm i Z_i)/Q$, with integers $A_i,B_i,Z_i$, and set
$C_i=B_i^2+Z_i^2\geq0$. At a grid shift $x=M/Q$, compute
\begin{align}
 D_1&=A_1-M,\label{eq:bit-pivots}\\
 D_i&=A_i-M-\left\lfloor\frac{C_{i-1}}{D_{i-1}}\right\rfloor
       \qquad(2\leq i\leq n),\notag
\end{align}
replacing every zero pivot by $-1$ immediately after computing it.
Thus every subsequent division has a nonzero integer denominator.
The floor is the signed mathematical floor, also for negative
denominators. Let $\nu$ be the number of negative pivots.

\begin{lemma}\label{lem:bit-count}
There is a real diagonal matrix $E$ with $\norm{E}\leq\eta$ such that
$\nu$ is the number of eigenvalues of $A+E$ strictly below $x$, and
$x$ is not an eigenvalue of $A+E$. Consequently, for every
$j\in\{1,\ldots,n\}$,
\begin{equation}\label{eq:bit-safe-decisions}
 \nu\geq j\ \Longrightarrow\ \lambda_j(A)<x+\eta,
 \qquad
 \nu<j\ \Longrightarrow\ \lambda_j(A)>x-\eta.
\end{equation}
\end{lemma}

\begin{proof}
Put $d_i=D_i/Q$. Before a zero replacement, the diagonal perturbation
relative to the exact scalar $LDL^*$ recurrence is
\[
 e_i=\eta\left(
       \frac{C_{i-1}}{D_{i-1}}-
       \left\lfloor\frac{C_{i-1}}{D_{i-1}}\right\rfloor
       \right)\in[0,\eta)\qquad(i\geq2).
\]
Replacing a zero pivot by $-1$ subtracts $\eta$ from this error,
leaving it in $[-\eta,0)$. The first error is either $0$ or
$-\eta$. Hence the computed $d_i$ are the exact $LDL^*$ pivots of
$A+\diag(e_i)-x\Id$. Its unit lower bidiagonal factor is
nonsingular, and the off-diagonal entries are unchanged. Congruence
therefore identifies its inertia with the signs of the nonzero
$d_i$. Their product also proves nonsingularity. This proves the
first assertion, including zero edges, repeated eigenvalues of $A$
and complex edge entries. Weyl's inequality and $\norm{E}\leq\eta$
give~\eqref{eq:bit-safe-decisions}.
\end{proof}

The auxiliary $E$ may depend on the query $x$. Both decisions in
\eqref{eq:bit-safe-decisions} concern the same original $A$, so the
construction does not need these auxiliary perturbations to agree
across queries.

\subsection{Fixed-grid spectral enclosures and word sizes}
\label{subsec:bit-bisection}

Suppose $\norm{A}<2$. For each $j=1,\ldots,n$, initialize integer
endpoints $L=-2Q$ and $U=2Q$. Perform exactly $p$ steps:
\[
 M=\left\lfloor\frac{L+U}{2}\right\rfloor,
 \qquad\text{compute $\nu$ by~\eqref{eq:bit-pivots}},
\]
and update
\begin{equation}\label{eq:bit-bisection}
 \begin{cases}
 U\gets\min(U,M+1),&\nu\geq j,\\
 L\gets\max(L,M-1),&\nu<j.
 \end{cases}
\end{equation}
The strict bracket $L/Q<\lambda_j(A)<U/Q$ is preserved by
Lemma~\ref{lem:bit-count}. If $w=U-L$, its next value is at most
$\lceil w/2\rceil+1\leq w/2+3/2$. Since $w_0=4Q$, after $p$
steps,
\begin{equation}\label{eq:bit-width}
 w_p\leq4+3(1-2^{-p})<7.
\end{equation}
The output interval thus has width at most $7\eta$, and its
midpoint has error at most $7\eta/2$. Intervals retain their ordered
index $j$; no sorting of their midpoints is needed, even with
repeated eigenvalues.

Each inertia count uses $O(n)$ integer operations, so all $n$
brackets cost $O(n^2p)$ such operations. In the applications below,
\begin{equation}\label{eq:bit-entry-bounds}
 |A_i|\leq2Q,\qquad |B_i|,|Z_i|\leq Q+1,
 \qquad |M|\leq2Q.
\end{equation}
Thus $C_i\leq8Q^2$. Since every denominator has absolute value at
least one,
\[
 \left|\left\lfloor C_i/D_i\right\rfloor\right|
       \leq8Q^2+1,\qquad |D_i|\leq13Q^2.
\]
All recurrence integers have at most $2p+O(1)$ bits. No rational
numerator or denominator is propagated through this recurrence.
Addition, multiplication and signed floor division cost $O(p^2)$
bit operations by elementary algorithms. The spectral-enclosure
oracle therefore has cost $O(n^2p^3)$ in bit operations and
$O(np)$ storage.

\subsection{One precision for every adaptive value query}
\label{subsec:bit-values}

The componentwise input enclosure implies a uniform norm bound.
Each complex edge error has magnitude at most $2\eta$, so the
Hermitian error matrix has absolute row sums at most $5\eta$.
For every admissible $T$ represented by the same proxy,
\begin{equation}\label{eq:bit-enclosure-norm}
 \norm{T-T_0}\leq5\eta.
\end{equation}
Use the convex function $F(t)=f(T+tX)$ of
\eqref{eq:X-W-F}, and start the two-child selector at $[0,\beta]$.
All queried parameters belong to this interval and have dyadic
denominators dividing $2^{a+2K+2}$, which divides $Q$ by
\eqref{eq:bit-precision}. The matrices $T_0+tX$ are therefore
represented exactly on the fixed grid. They satisfy
\[
 \norm{T_0+tX}\leq1+5\eta+n\beta<2,
 \qquad n\beta\leq\delta/4\leq1/16,
\]
and their entries obey~\eqref{eq:bit-entry-bounds}. Here $p\geq20$;
the edge bounds also follow directly from the component enclosure
and $\norm{T}\leq1$.

At each $t$, form the $n$ spectral brackets from
Section~\ref{subsec:bit-bisection}, and let $\widehat F(t)$ be the
sum of their midpoints weighted by $j$. Its error against
$f(T_0+tX)$ is at most $(7/2)W\eta$, where $W=n(n+1)/2$.
The Lipschitz bound~\eqref{eq:lipschitz} and
\eqref{eq:bit-enclosure-norm} imply, simultaneously for every
admissible $T$,
\begin{equation}\label{eq:bit-value-error}
 |\widehat F(t)-F(t)|
 \leq\frac{17}{2}W\eta
 \leq9n^2\eta
 \leq\frac{9\ell}{2^{20}}
 <\frac{\ell'}{1024}
 \qquad(\ell'\geq\ell).
\end{equation}
We used $n^2\eta\leq\ell/2^{20}$, which follows from the
definition of $m$. In particular this error allowance holds for
every child length during all $K$ rounds, with a single precision.

Apply Lemma~\ref{lem:robust-choice}, always choosing the smaller
computed child potential and resolving ties to the left. The two
children have the same length, so comparing their potentials amounts
to comparing the integer numerators of their four-value differences:
all weighted midpoint sums have the common denominator $2Q$.
The numerators have $O(p+\log n)=O(p)$ bits. These decisions compare
explicit finite integers, not the signs of unknown exact spectral
quantities.

The proof of~\eqref{eq:final-potential} applies to the initial
interval $[0,\beta]$, since its potential is still less than $n^3$.
The selected interval has length $\ell$ and center $c$. With
$r=\ell/4$, Lemma~\ref{lem:children} gives
\begin{equation}\label{eq:bit-small-defect}
 \Delta_r F(c):=F(c+r)+F(c-r)-2F(c)<r/8.
\end{equation}
Every decision depends only on $T_0$, and the error guarantee
\eqref{eq:bit-value-error} holds for the entire enclosure. Hence
the same $c$ satisfies~\eqref{eq:bit-small-defect} for every
admissible $T$.

The actual matrix $H=T+cX$ has norm at most $1+\delta/4<2$,
unchanged edge magnitudes at most one, and $0<r<1$. After the
diagonal phase gauge of Remark~\ref{rem:gauge} when necessary,
the contrapositive of Lemma~\ref{lem:defect} yields
$\gap(T+cX)>r^4/(10^{12}n^{14})$. Minimality in
\eqref{eq:bit-a}--\eqref{eq:bit-K} gives
$\beta>\delta/(8n)$ and $2^K<32n^3$, so
\begin{equation}\label{eq:bit-gap-constant}
 r=\frac{\beta}{4\cdot4^K}>
       \frac{\delta}{2^{15}n^7},\qquad
 \frac{r^4}{10^{12}n^{14}}>
       \frac{\delta^4}{10^{12}2^{60}n^{42}}>
       2^{-100}\frac{\delta^4}{n^{42}}.
\end{equation}
The last strict comparison uses $10^{12}<2^{40}$. The budget
follows from $0\leq c\leq\beta$ and $n\beta\leq\delta/4$.
This proves~\eqref{eq:bit-guarantee}.

\subsection{An independently checked certificate for the enclosure}
\label{subsec:bit-certificate}

Run the spectral-enclosure oracle at the three parameters
$t_-=c-r$, $t_0=c$ and $t_+=c+r$. Retain brackets
$[l_j(t),u_j(t)]$ of width at most $7\eta$ for each increasingly
ordered eigenvalue of $T_0+tX$, and define
\begin{align}
 L(t)&=\sum_{j=1}^n j\,l_j(t),&
 U(t)&=\sum_{j=1}^n j\,u_j(t),\notag\\
 B&=U(t_+)+U(t_-)-2L(t_0)+20W\eta.
 \label{eq:bit-certificate-bound}
\end{align}
By~\eqref{eq:bit-enclosure-norm}, each $F(t)$ lies between
$L(t)-5W\eta$ and $U(t)+5W\eta$. Therefore
$\Delta_rF(c)\leq B$ for every admissible $T$. Each of these
value intervals has length at most $17W\eta$, whence
\[
 0\leq B-\Delta_rF(c)\leq68W\eta.
\]
Using $W\leq n^2$, $n^2\eta\leq4r/2^{20}$ and
\eqref{eq:bit-small-defect}, we obtain the strict slack
\begin{equation}\label{eq:bit-certificate-slack}
 B<\frac r8+\frac{68r}{2^{18}}<\frac r4<\frac r2.
\end{equation}
The certificate verifier checks the parameters and budget, the three
sets of ordered spectral enclosures, and the rational inequality
$B<r/2$. It recomputes $B$ from those data. Equations
\eqref{eq:defect-lower}, \eqref{eq:bit-enclosure-norm} and
\eqref{eq:bit-certificate-bound} then prove the same gap bound
uniformly throughout the normalized part of the enclosure.

For completeness, the brackets can be verified independently by
exact rational inertia, without importing the safeguarded constructor.
To count negative eigenvalues of the Hermitian tridiagonal matrix
$A-x\Id$, eliminate a nonzero first pivot by its scalar Schur
complement and add its sign to the inertia. If the first pivot is
zero and the next edge is zero, remove that isolated zero. If the
first pivot is zero and the next edge is nonzero, use the leading
$2\times2$ block. Its determinant is $-|b|^2<0$, so it contributes
one positive and one negative eigenvalue. The $(2,2)$ entry of its
inverse is zero, so its Schur complement leaves the next diagonal
unchanged. Repeating these rules gives the exact count $N_<(x)$ of
eigenvalues strictly below $x$, including eigenvalue endpoints and
zero edges. All operations are rational, also for complex rational
entries, through their real and imaginary components.

For a claimed $j$-th bracket, the verifier requires
\begin{equation}\label{eq:bit-verifier-brackets}
 N_<(l_j)<j,\qquad N_<(u_j)\geq j.
\end{equation}
These imply $l_j\leq\lambda_j(A)<u_j$. Thus the verified lower
endpoint is inclusive and the upper endpoint is strict. The brackets
constructed earlier satisfy the stronger strict lower inequality as
well. No numerical eigenvalue routine enters this verifier.

There are three pairs of $n$ grid endpoints and a bounded number of
parameters; with the supplied proxy, their total size is $O(np)$
bits. The verifier receives the intended proxy separately and checks
its agreement with the certificate data. A passing certificate is a
uniform conditional guarantee for every normalized matrix in that
enclosure. It does not certify that such a matrix exists. Exact
rational verification can create larger intermediate integers, and
no bounded-word cost for this independent verifier is asserted.

\subsection{Total cost, supporting checks and scope}
\label{subsec:bit-cost-scope}

There are at most $8K$ value calls during selection and three further
calls for certification. Each costs $O(n^2p)$ operations on
$O(p)$-bit integers. Parameter construction, edge squaring, weighted
summation and the $n$ diagonal output entries fit within the same
bound. Since $a=O(d+\log n)$ and $K,m=O(\log n)$,
$p=O(\log(n/\delta))$. This proves~\eqref{eq:bit-cost} and its
schoolbook bit-operation bound. The input, the current interval,
one set of spectral brackets and the final certificate require
$O(np)$ workspace. These estimates are uniform in $n,d$ and the
input components, completing the proof of
Theorem~\ref{thm:finite-bit}.

The accompanying integer constructor and separate rational verifier
have exact supporting fixtures of orders $2,3,4$ and $8$.
They include complex edges, repeated entries, zero and weak edges,
large real couplings, and a dyadic proxy for a nondyadic rational
matrix. The rational source matrices' norm and enclosure promises
are checked exactly. Additional cases exercise exact inertia,
safeguarded-count branches and large-integer serialization.
Malformed certificates, altered brackets, incorrect defect data and
input mismatches are required to fail. The registered outputs,
certificate bytes and reason-specific rejection messages agree in
ordinary, optimized and alternate hash-seed executions. These finite
checks support the implementations; the uniform gap, word-size,
certificate-size and runtime bounds follow from the written proof,
not from the tested dimensions.

The result is a nearly quadratic finite-bit baseline. It does not
supply the missing nearly linear finite-precision spectral-value
oracle, nor provide deterministic regularization for general nonnormal
eigenproblems. The exact-arithmetic Graeffe evaluator
in Section~\ref{sec:evaluation} cannot simply be assigned the fixed
precision used here: its coefficient transformations can have large
cancellation, unlike the safeguarded recurrence
\eqref{eq:bit-pivots}. No numerical-stability theorem for a
floating-point implementation is claimed.

Dyadic budgets lose no qualitative generality. For an arbitrary
requested real budget $0<\eps<1/2$, a supplied dyadic $\delta$ with
$\eps/2\leq\delta\leq\eps$ and $\delta\leq1/4$ gives the same
conclusion with only an absolute-factor change in the gap constant.
This formulation still requires the finite proxy at the precision
determined by that supplied dyadic budget; it does not extract bits
from an arbitrary real input.
